\section{Introduction}

% Recent advancements in imitation learning architecture have pushed the boundaries of robotic autonomy in complex environments~\citep{ai2025review}. Diffusion-based policies~\citep{ddpm, flowmatching, dp, intelligence2025pi}, as one such breakthrough, effectively capture the inherent multi-modality of human behavior by formulating action prediction as a conditional denoising process. In this paradigm, a neural network is trained to estimate injected noise during training, while at inference time, actions are sampled from random noise and gradually denoised into executable control commands.

Recent advances in imitation learning architectures have significantly expanded the capabilities of robotic systems operating in complex and unstructured environments~\citep{ai2025review}. Among these, diffusion-based policies~\citep{ddpm, flowmatching, dp, intelligence2025pi} have emerged as a powerful paradigm for modeling the intrinsic multi-modality of human demonstrations. These methods formulate action generation as a conditional denoising process: during training, a neural network learns to predict injected noise, while at inference time, executable actions are obtained by iteratively denoising samples initialized from random noise.

% Diffusion models~\citep{ddpm,song2020denoising,rombach2022high,flowmatching,rectifiedflow} have revolutionized generative modeling, demonstrating remarkable capabilities in high-fidelity image synthesis and video generation. 
% Building on this success, recent robotics research has recast action prediction as a conditional denoising process~\citep{dp, intelligence2025pi}. 


Despite their strong empirical performance on high-precision and multi-modal tasks, diffusion models suffer from a well-known limitation: the ``denoise-from-scratch'' paradigm incurs substantial inference latency~\citep{pan2025much}. Generating a single action typically requires dozens of iterative denoising steps, creating a major bottleneck for real-time robotic control, where low cycle time and rapid feedback are essential for stable execution.
To mitigate this limitation, prior work has explored improving the initialization of the diffusion process to accelerate inference. For instance,~\citet{steeringdp} propose an RL-trained policy to steer the initial sampling distribution, while~\citet{warmstarts} employ warm-start strategies to identify more informative starting points. These approaches replace uninformed \textit{Gaussian} noise with priors closer to the data distribution, often through auxiliary models that predict distributional statistics conditioned on the current observation. While effective, such approaches still rely on stochastic noise initialization and inevitably introduce additional modeling complexity.

This observation raises a more fundamental question: \textbf{do robotic policies truly need to be generated by sampling from random noise?} Diffusion models were originally developed for high-fidelity image synthesis and video generation~\citep{ddpm,song2020denoising,rombach2022high,flowmatching,rectifiedflow,blattmann2023stable}, where generation typically begins from uninformed noise due to the absence of meaningful priors. Robot control, however, operates under a fundamentally different regime. Modern robots are equipped with rich state sensors that provide continuous, low-latency feedback about the system’s configuration and dynamics~\citep{jia2024feedback, jia2024learning}. This structured feedback constitutes a strong and reliable prior, naturally encoding the robot’s current physical state or recent execution history, and thus offers a principled alternative to random noise initialization for action generation~\citep{jia2025foreseer}.

% Modern robotic systems are equipped with sensors that provide continuous and precise proprioceptive action/state feedback, offering a rich and reliable source of prior information~\citep{jia2024feedback}. Such information naturally captures the robot’s physical state and dynamics, and can serve as a strong initialization signal for action generation~\citep{jia2025foreseer}.
% potentially eliminating the need for uninformed noise sampling altogether. 

However, as shown in Fig.~\ref{framework} (a) and (b), most existing approaches adhere to either regression- or diffusion-based paradigms that condition action generation on a simple concatenation of the current proprioceptive state/action and visual observations. Such designs overlook the temporal continuity intrinsic to robotic motion and often dilute low-dimensional proprioceptive signals when fused directly with high-dimensional visual representations. Moreover, recent studies show that explicit conditioning on proprioceptive states can adversely affect spatial generalization~\citep{zhao2025you}.
Consequently, rich historical information about system dynamics and action trends remains largely underexploited in prevailing frameworks. Notably, diffusion models define mappings between probability distributions rather than individual states. This perspective motivates a natural question: can the inherent proximity between distributions of past executions and future actions be exploited to reduce learning complexity, yielding a shorter and more stable transport path for policy generation?

% 缺一块讨论：这里要更加深入的讨论，Motivated by xx，我们是如何从理论上考虑的，比如本质是分布转换，那么利用本体感知的历史到xx的转换是否可以减少难度，这一简单的改变是否就可以带来好的特性？因为理论上是更加成立的？这一段是一种假设，并让读者读懂你“为什么有了这个idea”。
% 下一段是To this end, we propose xx

To this end, we propose \textbf{A}ction-\textbf{to}-\textbf{A}ction {Flow} Matching (\textbf{\m{}}), a novel flowing matching-based policy paradigm that shifts action generation \textit{from uninformed sampling to informed initialization}. Unlike prior approaches that initialize the diffusion process with standard \textit{Gaussian} noise, \m{} directly leverages a sequence of historical proprioceptive actions as the starting point for action generation, as illustrated in Fig.~\ref{framework}. To capture subtle motion patterns and temporal dependencies, these low-dimensional action histories are embedded into a high-dimensional latent space. By learning a flow that transports historical action distributions to future actions, \m{} bypasses the costly iterative denoising process from \textit{Gaussian} noise.
We evaluate \m{} extensively in both simulated environments and real-world robotic systems. Empirically, our method exhibits remarkable training efficiency and consistently outperforms 8 state-of-the-art baselines.
In particular, \m{} achieves significantly faster inference convergence, i.e., up to 20$\times$ and $5\times$ faster than vanilla diffusion and flow matching methods, respectively, while \textbf{enabling high-quality action generation in as few as a single inference step}. 
Moreover, grounding the generation process in proprioceptive history substantially improves robustness to visual perturbations, and the history-informed initialization enhances generalization to unseen configurations by enforcing physical consistency over time.

In summary, we introduce a new diffusion-based robot policy paradigm that replaces stochastic noise initialization with history-based proprioceptive initialization, enabling informed action generation grounded in the robot’s own dynamics. Leveraging latent state representations, \m{} captures fine-grained motion structure and supports efficient action-to-action transitions without iterative denoising from random noise. Extensive experiments in both simulated and real-world robotic environments demonstrate that our method achieves state-of-the-art performance across training efficiency, inference speed, robustness to visual perturbations, and generalization to unseen configurations. Furthermore, we showcase the applicability of \m{} in robotic video generation, suggesting promising potential for broader scalability.